\documentclass[../DoAn.tex]{subfiles}
\begin{document}

\section{Đặt vấn đề}
\label{section:1.1}

Giáo dục hiện đại đòi hỏi sự hỗ trợ học tập cá nhân hóa và tự động hóa các tác vụ giáo dục phức hợp. Các hệ thống truyền thống như tìm kiếm thông tin hoặc câu hỏi-trả lời đơn giản không thể vượt qua hiện tượng "ảo giác" (hallucination) của các mô hình ngôn ngữ lớn và thiếu độ chính xác trong ngữ cảnh giáo dục.

Ở cấp bậc THPT, nhu cầu hỗ trợ học tập ngày càng cấp thiết. Học sinh cần một trợ lý ảo có thể trả lời câu hỏi dựa trên tài liệu chính thức như sách giáo khoa, tự động sinh ra bài tập với mức độ tùy chỉnh, chấm điểm và giải thích lỗi một cách chính xác. Đồng thời, giáo viên phải đầu tư nhiều thời gian cho việc tạo bài tập, chấm điểm, và giảng dạy lặp lại nội dung. Các hệ thống AI chung ngoài thị trường thường không tuân thủ chương trình giáo dục Việt Nam và dễ sinh ra thông tin sai lệch, làm giảm hiệu quả dạy-học.

Vấn đề cần giải quyết là xây dựng một hệ thống trợ lý ảo thông minh có khả năng tích hợp nhiều chức năng giáo dục (trả lời câu hỏi, sinh đề bài, chấm điểm, tạo tài liệu) mà vẫn đảm bảo tính chính xác cao dựa trên sách giáo khoa, giảm bớt gánh nặng cho giáo viên và nâng cao chất lượng học tập cho học sinh tại các trường THPT.

\section{Mục tiêu và phạm vi đề tài}
\label{section:1.2}

Hiện nay, các giải pháp hỗ trợ giáo dục AI đã tồn tại trên thị trường. Các hệ thống như ChatGPT hoặc Gemini cung cấp khả năng tương tác linh hoạt nhưng dễ sinh ra thông tin không chính xác khi áp dụng cho bối cảnh giải pháp cụ thể. Các hệ thống quản lý lớp học truyền thống (Moodle, Google Classroom) cung cấp công cụ quản lý tập trung nhưng thiếu khả năng AI tự động. Các nền tảng hỗ trợ giáo dục AI chuyên biệt từ các nước ngoài thường xây dựng cho nền tảng giáo dục quốc tế, không tối ưu cho chương trình giáo dục Việt Nam và yêu cầu chi phí cao.

Các hạn chế của các giải pháp hiện có là thiếu tính "gắn kết dữ liệu" (grounding) với tài liệu chính thức, dẫn tới sai lệch nội dung. Hầu hết các hệ thống chỉ hỗ trợ một hoặc hai chức năng đơn lẻ, không cung cấp quy trình giảng dạy toàn diện. Đặc biệt, không có cơ chế tự kiểm tra (self-reflection) chất lượng của phản hồi, dẫn tới độ tin cậy thấp.

Mục tiêu chính của đồ án là xây dựng một hệ thống tên gọi "Intelligent Educational Assistant" (Trợ lý Ảo Thông Minh Hỗ Trợ Giáo Dục) với bốn chức năng cốt lõi: (i) trả lời câu hỏi (Q\&A) dựa trên sách giáo khoa Cánh Diều và Kết Nối Tri Thức (lớp 10, 11, 12), (ii) sinh tự động đề bài tập với nhiều định dạng (trắc nghiệm, điền khuyết, tự luận, đúng/sai) có độ khó tùy chỉnh, (iii) chấm điểm tự động và cung cấp giải thích lỗi dựa trên ngữ cảnh chính thức, (iv) sinh tài liệu học tập như bài giảng hoặc bản tóm tắt từ nội dung sách giáo khoa.

Phạm vi của đồ án bao gồm xử lý dữ liệu từ hai bộ sách giáo khoa (Cánh Diều và Kết Nối Tri Thức) môn Tin học lớp 10, 11, 12, số lượng tương ứng là 2.348 đoạn dữ liệu sau khi xử lý theo cấu trúc phân tầng. Hệ thống tập trung vào bốn chức năng chính đã nêu, với đối tượng người dùng là học sinh và giáo viên bậc THPT.

\section{Định hướng giải pháp}
\label{section:1.3}

Hệ thống sử dụng kiến trúc Retrieval-Augmented Generation (RAG) kết hợp với Multi-Agent Orchestration để giải quyết các thách thức đã nêu. RAG là một phương pháp giảm hiện tượng ảo giác bằng cách truy xuất các đoạn dữ liệu liên quan từ kho tài liệu (sách giáo khoa), sau đó cho mô hình ngôn ngữ tổng hợp phản hồi dựa trên ngữ cảnh chính thức này.

Giải pháp được xây dựng từ ba thành phần chính. Thứ nhất, thành phần truy xuất thông tin (Retrieval) sử dụng "Hybrid Search" kết hợp tìm kiếm từ khóa (Custom BM25 viết từ đầu) và tìm kiếm ngữ nghĩa (Cosine Similarity trên embedding vector). Thứ hai, thành phần tổng hợp phản hồi (Generation) sử dụng Google Gemini 2.5 làm mô hình ngôn ngữ chính, kết hợp với các "Specialist Handlers" để xử lý từng loại tác vụ cụ thể. Thứ ba, thành phần tự kiểm tra (Self-Reflection) chạy một bộ hợp thức (Validator Agent) để đảm bảo phản hồi tuân thủ ngữ cảnh gốc và đáp ứng yêu cầu.

Lý do lựa chọn RAG là phương pháp này đã chứng minh hiệu quả trong việc giảm ảo giác và nâng cao độ tin cậy của mô hình ngôn ngữ. Custom BM25 được lựa chọn vì kho dữ liệu nhỏ (2.348 đoạn), không cần công cụ tối ưu hóa phức tạp như FAISS, và dễ giải thích trong báo cáo. Việc sử dụng cả tìm kiếm từ khóa và ngữ nghĩa đảm bảo phủ phát toàn diện các loại truy vấn khác nhau của người dùng. Google Gemini 2.5 được chọn vì hiệu năng cao, API ổn định, và chi phí hợp lý so với các mô hình khác.

Các đóng góp chính của đồ án bao gồm: (i) xây dựng một hệ thống giáo dục tích hợp đầu tiên kết hợp bốn chức năng thành một nền tảng duy nhất, (ii) thiết kế một pipeline RAG thích ứng (Adaptive RAG) lựa chọn động chiến lược truy xuất dựa trên loại truy vấn, (iii) kiến trúc Multi-Agent Orchestration linh hoạt cho phép mở rộng chức năng dễ dàng mà không ảnh hưởng tới các thành phần hiện có, (iv) cô lập hiện tượng ảo giác của mô hình thông qua cơ chế tự kiểm tra (Self-Reflection), (v) đo lường định lượng chất lượng hệ thống sử dụng RAGAS evaluation framework với bốn chỉ số (Answer Relevancy, Faithfulness, Context Precision, Context Recall).

\section{Bố cục đồ án}
\label{section:1.4}

Báo cáo đồ án tốt nghiệp này được tổ chức thành sáu chương chính như sau.

Chương 2 trình bày các nền tảng lý thuyết cơ bản hỗ trợ cho việc thiết kế hệ thống. Các khái niệm Retrieval-Augmented Generation, Multi-Agent Architecture, các mô hình embedding tiếng Việt, và các kỹ thuật xử lý ngôn ngữ tự nhiên liên quan được giới thiệu chi tiết. Chương này tạo nên cơ sở lý thuyết vững chắc cho các quyết định thiết kế trong các chương tiếp theo.

Chương 3 mô tả chi tiết các công nghệ và phương pháp được sử dụng trong hệ thống. Các thành phần kỹ thuật được trình bày bao gồm Custom BM25 search engine viết từ đầu, tìm kiếm ngữ nghĩa sử dụng mô hình embedding tiếng Việt, tìm kiếm kết hợp (Hybrid Search) với bình thường hóa Reciprocal Rank Fusion (RRF), tái xếp hạng kết quả (Cross-Encoder Reranking), và chiến lược phân chia tài liệu theo cấu trúc phân tầng (Hierarchical Chunking). Mỗi phương pháp được phân tích về lợi ích, hạn chế, và lý do lựa chọn.

Chương 4 trình bày kiến trúc tổng thể của hệ thống và quá trình triển khai. Các thành phần chính bao gồm Intent Detection pipeline để phân loại ý định người dùng, Adaptive RAG Agent để lựa chọn chiến lược truy xuất phù hợp, Specialist Handlers cluster để xử lý từng tác vụ cụ thể (Chat, Quiz Generation, Answer Scoring, Content Generation), và cơ chế Self-Reflection đảm bảo chất lượng. Phần này cũng mô tả cách các module tương tác với nhau thông qua một pipeline xuyên suốt.

Chương 5 trình bày các đóng góp chính của đồ án cùng với kết quả thực nghiệm. Phần này bao gồm mô tả giao diện người dùng được xây dựng bằng Gradio, các kết quả định lượng sử dụng RAGAS framework, và các trường hợp nghiên cứu minh họa khả năng của hệ thống trong các tình huống sử dụng thực tế.

Chương 6 nêu rõ các hạn chế hiện tại của hệ thống như độ chính xác của Intent Detection, tốc độ xử lý, và quy mô dữ liệu. Báo cáo đề xuất các hướng cải tiến trong tương lai bao gồm fine-tuning mô hình Intent Detection, triển khai cơ chế caching để tăng tốc độ, và mở rộng hệ thống sang các môn học khác. Cuối cùng, chương này cung cấp các nhận xét kết luận về tác dụng học tập của hệ thống và triển vọng phát triển lâu dài.

\end{document}
